\documentclass{beamer}
%Information to be included in the title page:
\title{RL-LDPC presentation}
\author{}
\institute{}
\date{11th April 2026}

\begin{document}

\frame{\titlepage}
% -------------------------------------------------------------------------

\begin{frame}{Outline}
  \tableofcontents
\end{frame}

\section{Mutual Information}
% Contents Slide
\begin{frame}{Overview}
\begin{itemize}
    \item Implemented naive mutual information based decoding
    \item Changed RELDEC MDRP formulation based on Mutual information
\end{itemize}
\end{frame}
% -------------------------------------------------------------------------
\subsection{Intro to Mutual Info}
\begin{frame}{Mutual Information: Intuition}
\begin{itemize}
    \item It quantifies how much information messages carry about transmitted bits
    \item Let $X \in \{0,1\}$ be the transmitted bit
    \item Let $L$ be the LLR message
    \item Mutual Information $I(X;L)$ measures:
    \begin{block}{}
        Reduction in uncertainty about $X$ after observing $L$
    \end{block}
    \item Range:
    \begin{itemize}
        \item $0$ $\rightarrow$ no information
        \item $1$ $\rightarrow$ perfect knowledge
    \end{itemize}
\end{itemize}
\end{frame}

%------------------------------------------------
\begin{frame}{Formal Definition}
\begin{itemize}
    \item Mutual Information:
\end{itemize}

\[
I(X;L) = \mathbb{E}\left[\log_2 \frac{p(L|X)}{p(L)}\right]
\]

\begin{itemize}
    \item $X \in \{0,1\}$, $L$ is the LLR
    \item Typically assume all-zero codeword (symmetry)
\end{itemize}
\end{frame}

%------------------------------------------------
\begin{frame}{MI and LLRs}
\begin{itemize}
    \item For symmetric channels, MI depends only on LLR distribution
    \item Practical computation:
    \begin{itemize}
        \item Generate LLR samples
        \item Estimate MI empirically
    \end{itemize}
\end{itemize}
\end{frame}

%------------------------------------------------
\begin{frame}{Mutual Information for a Cluster}
\begin{itemize}
    \item Consider a cluster of check nodes: $\mathcal{C}_k$
    \item Collect all outgoing LLR messages from this cluster:
    \[
    \{L_i : i \in \mathcal{C}_k\}
    \]

    \item Treat these messages as samples from a random variable $L^{(k)}$

    \item Define cluster MI as:
    \[
    I_k = I(X; L^{(k)})
    \]

    \item Interpretation:
    \begin{itemize}
        \item Measures average information content of messages from cluster $\mathcal{C}_k$
        \item Each cluster has its own MI evolution
    \end{itemize}
\end{itemize}
\end{frame}
%------------------------------------------------
%------------------------------------------------
\begin{frame}{Computing MI}
\begin{itemize}
    \item Assume LLRs are Gaussian:
    \[
    L \sim \mathcal{N}(\mu, 2\mu)
    \] 
    \item Exact formula (when conditioned on $x_j = +1$), we have
$$I_{E,V} = 1 - \int_{-\infty}^{\infty} \frac{1}{\sqrt{2\pi}\sigma} e^{-(l-\sigma^2/2)^2/(2\sigma^2)} \log_2(1+e^{-l}) dl. $$

For convenience we write this as
$$I_{E,V} = J(\sigma) = J\left(\sqrt{(d_v-1)\sigma_A^2 + \sigma_{ch}^2}\right). $$
        \item Approximation: 
For computer implementation, the $J(\sigma)$ function is split into two intervals based on a threshold of $\sigma^* = 1.6363$:

$$J(\sigma) \approx \begin{cases} 
a_{J,1} \sigma^3 + b_{J,1} \sigma^2 + c_{J,1} \sigma, & 0 \leq \sigma \leq \sigma^* \\
1 - e^{a_{J,2} \sigma^3 + b_{J,2} \sigma^2 + c_{J,2} \sigma + d_{J,2}}, & \sigma^* < \sigma < 10 \\
1, & \sigma > 10 
\end{cases}$$

\end{itemize}
\end{frame}

%------------------------------------------------

%------------------------------------------------
\begin{frame}{Apprximation}
\begin{itemize}

        \item Approximation: 
For computer implementation, the $J(\sigma)$ function is split into two intervals based on a threshold of $\sigma^* = 1.6363$:

$$J(\sigma) \approx \begin{cases} 
a_{J,1} \sigma^3 + b_{J,1} \sigma^2 + c_{J,1} \sigma, & 0 \leq \sigma \leq \sigma^* \\
1 - e^{a_{J,2} \sigma^3 + b_{J,2} \sigma^2 + c_{J,2} \sigma + d_{J,2}}, & \sigma^* < \sigma < 10 \\
1, & \sigma > 10 
\end{cases}$$

Numerical Constants for $J(\sigma)$:
\begin{itemize}
    \item $a_{J,1} = -0.0421061$
    \item $b_{J,1} = 0.209252$
    \item $c_{J,1} = -0.00640081$
    \item $a_{J,2} = 0.00181491$
    \item $b_{J,2} = -0.142675$
    \item $c_{J,2} = -0.0822054$
    \item $d_{J,2} = 0.0549608$
\end{itemize}
\end{itemize}
\end{frame}

%------------------------------------------------

\begin{frame}{MI Evolution in Iterations}
\begin{itemize}
    \item Early iterations:
    \begin{itemize}
        \item Low MI (noisy information)
    \end{itemize}
    \item As iterations proceed:
    \begin{itemize}
        \item Messages become more reliable
        \item MI increases
    \end{itemize}
    \item Outcomes:
    \begin{itemize}
        \item Successful decoding: $I \rightarrow 1$
        \item Failure: MI saturates below 1
    \end{itemize}
\end{itemize}
\end{frame}

%------------------------------------------------

%------------------------------------------------

\begin{frame}{EXIT Chart (High-Level)}
\begin{itemize}
    \item Plot relationship between input and output MI
    \item Axes:
    \begin{itemize}
        \item x-axis: $I_A$ (output)
        \item y-axis: $I_E$ (input)
    \end{itemize}
    \item Two curves:
    \begin{itemize}
        \item Variable node curve
        \item Check node curve
    \end{itemize}
    \item Decoding succeeds if curves form an ``open tunnel''
\end{itemize}
\begin{figure}
    \centering
    \includegraphics[width=0.5\linewidth]{image9.png}
    \label{fig:placeholder}
\end{figure}
\end{frame}

%------------------------------------------------
\subsection{Algorithms and Results}

\begin{frame}{Naive MI based decoding}
\begin{itemize}
    \item Goal: select the next check-node cluster $\mathcal{C}_k$ that yields the largest information gain

    \item For each cluster, track input mutual information:
    \[
    I^{(k)}_{\text{in}} = I(X; L^{(k)})
    \]
    \item Compute expected gain from next iteration:
    \[
    \Delta I_k = I^{(k)}_{\text{out}} - I^{(k)}_{\text{in}}
    \]

    \item Select cluster with max $\Delta I_k\$

    \item Insight:
    \begin{itemize}
        \item Prioritize regions where decoding can improve most
        \item Avoid wasting updates on already reliable clusters
    \end{itemize}
\end{itemize}
\end{frame}

%------------------------------------------------

\begin{frame}
\frametitle{Mutual Information State Space Q-Learning (MI-Tabular)}
\begin{itemize}
    \item \textbf{State:} For cluster $c$, state is MI-binned:
    $$s_c = \left\lfloor I_c \cdot (B-1) \right\rfloor, \quad s_c \in \{0, 1, \ldots, B-1\}$$
    where $I_c$ is mutual information of LLR posteriors on cluster $c$ neighbors, $B$ is the number of bins (default 21).
    
    \item \textbf{Reward:} MI-gain from cluster update:
    $$r(a) = I_a^{\text{after}} - I_a^{\text{before}}$$    
    
    \item Algorithms: Epsilon greedy Q learning and DQN
    
    \item \textbf{Note: } State is continuous for DQN version
\end{itemize}


\end{frame}

%------------------------------------------------
\begin{frame}{Results}
\begin{figure}
    \centering
    \includegraphics[width=1\linewidth]{image5.png}
\end{figure}
\end{frame}


%------------------------------------------------
\begin{frame}{Results}
\begin{figure}
    \centering
    \includegraphics[width=1\linewidth]{image3.png}
\end{figure}
\end{frame}


\section{PPO+GNN method}

\section{Other MDP formulation changes}
\subsection{Directly using global state space}
\subsection{Note on Restless MABs and Whittles Index}
\subsection{Augmented State space}



%------------------------------------------------
\section{Ongoing tasks and Goals}

\begin{frame}{Ongoing task}
\begin{figure}
    \centering
    \includegraphics[width=1\linewidth]{WhatsApp Image 2026-04-22 at 23.05.32.jpeg}
    \caption{Deep RELDEC loss explodes after a point in training}
\end{figure}
\end{frame}

\begin{frame}{Immediate Goals}
\begin{itemize}
    \item Refactor (Reorganize) the codebase as it is not possible to work with it anymore
    \item Optimize training and evaluation
\end{itemize}
\end{frame}



\end{document}

